\documentclass[11pt,a4paper,oneside]{article}

% -----------------------------------------------------------------------------
% ImageProcessing internal architecture paper
% Self-contained source; no external bibliography database or shell escape.
% -----------------------------------------------------------------------------
\usepackage[T1]{fontenc}
\usepackage[utf8]{inputenc}
\usepackage{lmodern}
\usepackage{microtype}
\usepackage[a4paper,margin=25mm,headheight=15pt]{geometry}
\usepackage{amsmath,amssymb,mathtools}
\usepackage{booktabs,longtable,tabularx,array,multirow}
\usepackage{enumitem}
\usepackage{xcolor}
\usepackage{listings}
\usepackage{fancyhdr}
\usepackage{titlesec}
\usepackage{hyperref}
\usepackage[nameinlink,noabbrev]{cleveref}

\definecolor{ink}{HTML}{202A35}
\definecolor{muted}{HTML}{5D6875}
\definecolor{rulegray}{HTML}{D7DDE3}
\definecolor{codebg}{HTML}{F5F7F9}
\definecolor{codegreen}{HTML}{28623B}
\definecolor{codeblue}{HTML}{174A73}

\hypersetup{
  colorlinks=true,
  linkcolor=codeblue,
  citecolor=codegreen,
  urlcolor=codeblue,
  pdftitle={ImageProcessing: Architecture and Semantic Contracts},
  pdfauthor={ImageProcessing Project},
  pdfsubject={Internal software architecture design paper},
  pdfkeywords={image processing, operation objects, functional composition, compute providers, qualification, heterogeneous execution}
}

\setlist[itemize]{leftmargin=1.55em,itemsep=2pt,topsep=4pt}
\setlist[enumerate]{leftmargin=1.75em,itemsep=2pt,topsep=4pt}
\setlength{\parskip}{4pt plus 1pt minus 1pt}
\setlength{\parindent}{0pt}
\renewcommand{\arraystretch}{1.16}
\titleformat{\section}{\Large\bfseries\color{ink}}{\thesection}{0.7em}{}
\titleformat{\subsection}{\large\bfseries\color{ink}}{\thesubsection}{0.7em}{}
\titleformat{\subsubsection}{\normalsize\bfseries\color{ink}}{\thesubsubsection}{0.7em}{}
\pagestyle{fancy}
\fancyhf{}
\fancyhead[L]{\small\textcolor{muted}{ImageProcessing --- Internal Architecture Paper}}
\fancyhead[R]{\small\textcolor{muted}{Design baseline 0.3}}
\fancyfoot[L]{\small\textcolor{muted}{Internal engineering document --- not for publication}}
\fancyfoot[R]{\small\textcolor{muted}{\thepage}}
\renewcommand{\headrulewidth}{0.35pt}
\renewcommand{\footrulewidth}{0pt}

\lstdefinelanguage{SwiftLite}{
  morekeywords={public,private,internal,struct,class,enum,protocol,func,let,var,try,await,async,throws,return,where,associatedtype,extension,static,final,actor,if,else,guard,case,switch,for,in,nil,true,false,any,some},
  sensitive=true,
  morecomment=[l]{//},
  morestring=[b]"
}
\lstset{
  language=SwiftLite,
  basicstyle=\small\ttfamily,
  keywordstyle=\color{codeblue}\bfseries,
  commentstyle=\color{codegreen},
  stringstyle=\color{red!55!black},
  backgroundcolor=\color{codebg},
  frame=single,
  rulecolor=\color{rulegray},
  framerule=0.4pt,
  xleftmargin=0.7em,
  xrightmargin=0.7em,
  framexleftmargin=0.45em,
  framexrightmargin=0.45em,
  breaklines=true,
  columns=fullflexible,
  keepspaces=true,
  showstringspaces=false,
  tabsize=4
}

\newcommand{\Project}{\textsc{ImageProcessing}}
\newcommand{\Image}{\texttt{Image}}
\newcommand{\Layer}{\texttt{Layer}}
\newcommand{\PixelResource}{\texttt{PixelResource}}
\newcommand{\Operation}{\texttt{Operation}}
\newcommand{\Provider}{\texttt{ComputeProvider}}
\newcommand{\Impl}{\texttt{OperationImplementation}}
\newcommand{\Qualify}{\texttt{qualify}}
\newcommand{\Select}{\texttt{select}}
\newcommand{\FP}{functional programming}
\newcommand{\OOP}{object-oriented programming}
\newcolumntype{Y}{>{\raggedright\arraybackslash}X}
\newcolumntype{L}[1]{>{\raggedright\arraybackslash}p{#1}}
\newcolumntype{Z}[1]{>{\raggedright\arraybackslash}p{#1}}

\begin{document}

\begin{titlepage}
  \centering
  \vspace*{18mm}
  {\Large\bfseries IMAGEPROCESSING PROJECT\par}
  \vspace{13mm}
  {\huge\bfseries ImageProcessing\par}
  \vspace{3mm}
  {\LARGE Architecture and Semantic Contracts\par}
  \vspace{9mm}
  {\large A backend-neutral, operation-oriented design for layered images and heterogeneous compute\par}
  \vfill
  \begin{tabular}{@{}ll@{}}
    \textbf{Document type:} & Internal technical design paper \\
    \textbf{Status:} & Proposed architecture baseline; implementation not yet specified \\
    \textbf{Version:} & 0.3 \\
    \textbf{Date:} & 20 November 2026 \\
    \textbf{Repository:} & \url{https://github.com/Shennagyp/ImageProcessing} \\
    \textbf{Audience:} & Library architecture, API, runtime, backend and test engineering \\
  \end{tabular}
  \vfill
  \begin{minipage}{0.86\textwidth}
    \small\color{muted}
    This document is an engineering record and design instrument. It is deliberately not formatted as a publication, does not claim experimental validation, and does not describe functionality as implemented. It records the present architectural direction, the reasoning behind it, the constraints that must remain invariant, and the questions that must be resolved before implementation contracts are frozen.
  \end{minipage}
  \vspace*{13mm}
\end{titlepage}

\pagenumbering{roman}
\tableofcontents
\clearpage
\pagenumbering{arabic}

\section*{Document control}
\addcontentsline{toc}{section}{Document control}

\begin{longtable}{@{}L{0.22\textwidth}Z{0.72\textwidth}@{}}
\toprule
\textbf{Field} & \textbf{Definition} \\
\midrule
Purpose & Establish a coherent architecture for a reusable image-processing library with layered image objects, operation objects, compositional processing, and environment-aware resolution among heterogeneous implementations. \\
Scope & Public semantic model, object and resource model, operation/implementation specifications, composition, compute-provider protocols, environment observations, common qualification/resolution, execution, memory, numerical behavior, and verification strategy. \\
Out of scope & A node editor or visual programming frontend; a general workflow engine; a new machine-learning training framework; a replacement for OpenCV, PyTorch, or ONNX; and application-specific user interfaces. \\
Normative language & \textquotedblleft{}MUST\textquotedblright{}, \textquotedblleft{}MUST NOT\textquotedblright{}, \textquotedblleft{}SHOULD\textquotedblright{}, and \textquotedblleft{}MAY\textquotedblright{} express design strength within this proposal. They do not imply that code already enforces the requirement. \\
Evidence status & Statements about third-party systems are referenced to their official documentation. Architectural recommendations are design conclusions, not claims that a cited library implements this project's proposed model. \\
Change policy & Changes to operation semantics, image interpretation, mutation guarantees, or provider contracts require an explicit design revision and corresponding tests. \\
Revision 0.2 & Corrects protocol-library layering: \texttt{ImageProcessing\allowbreak Protocol} is the primary domain contract and depends one-way on the generic, image-agnostic \texttt{GeneralExecution\allowbreak Protocol}. \\ 
Revision 0.3 & Establishes one generic operation-implementation contract, declarative operation and implementation specifications, environment-aware qualification and routing, fluent pipeline policy, construction-time conflict validation, and independent implementation packages. \\
\bottomrule
\end{longtable}

\section{Executive design}

\subsection{The central model}

\Project{} should be object-oriented in its representation of image documents, operation-oriented in its processing API, functional in the semantics of ordinary transformations, and heterogeneous in its execution implementation.

The principal distinction is between an operation and any concrete implementation of it. An operation such as \texttt{GaussianBlur} defines one stable meaning, parameter model, input/output contract, and numerical contract. Multiple independent implementation classes may implement that same operation: for example, \texttt{CPUGaussianBlur}, \texttt{CUDAGaussianBlur}, and \texttt{MetalGaussianBlur}. They are separate implementations of one operation, not separate operations and not methods added to a monolithic backend interface.

The domain contract defines one generic \texttt{ImageOperationImplementation} protocol, specialized by an associated operation type. Each implementation publishes a declarative \texttt{ImplementationSpecification}; it does not need to supply a bespoke qualification algorithm. The runtime may type-erase implementations internally for heterogeneous registration while retaining typed operation and implementation APIs at module boundaries.

Four sources of information participate in execution resolution:
\begin{enumerate}
  \item \textbf{Operation and operation specification:} what the operation means, which inputs and outputs it accepts, and which semantic/numerical guarantees are required or permitted.
  \item \textbf{Implementation specification:} what one concrete operation implementation supports, its algorithm-specific limits, precision behavior, provider requirements, and resource characteristics.
  \item \textbf{Execution environment:} which platforms, providers, devices, capabilities, and resources are actually available, plus the freshest reliable state observations such as queue pressure and available memory.
  \item \textbf{Pipeline qualification policy:} declarative requirements and preferences supplied through a fluent API on the composed function.
\end{enumerate}

A common, data-driven qualifying/resolution function evaluates these complementary inputs. It first excludes candidates that cannot satisfy hard semantic or technical requirements, then ranks the remaining candidates using the pipeline objectives and current execution state. The function can therefore route a busy GPU workload to a compatible CPU implementation when the pipeline permits it, or select a different implementation that satisfies a required FP32 computation mode. It must never select an implementation that violates the operation's semantic contract.

The protocol-library direction is deliberate. \texttt{ImageProcessing\allowbreak Protocol} is the primary, domain-facing contract library. It defines image-processing semantics, operations, implementation specifications, and declarative qualification-policy types. It consumes generic execution facts from the lower-level \texttt{GeneralExecution\allowbreak Protocol}. The latter remains image-agnostic: it describes providers, devices, resources, submission, completion, generic capabilities, and observable execution state, but it does not know what an image, layer, or blur operation is.

Concrete execution-provider packages and image-operation implementation packages are outside the core protocol/runtime dependency. A consuming application selects and registers the packages it needs. Adding an implementation must not require changing the operation definition, qualification algorithm, or core runtime.

\subsection{The intended execution relationship}

\begin{center}
\begin{tabularx}{0.96\textwidth}{@{}>{\centering\arraybackslash}X@{}}
\fbox{\parbox{0.88\textwidth}{\centering\textbf{Image document and target layer}\\[2pt] Logical identity, layer semantics, resource references}}\\[5pt]
$\Downarrow$\\[5pt]
\fbox{\parbox{0.88\textwidth}{\centering\textbf{Immutable operation or composed function}\\[2pt] Stable operation identity, parameters, operation specification, pipeline policy}}\\[5pt]
$\Downarrow$\\[5pt]
\fbox{\parbox{0.88\textwidth}{\centering\textbf{Common qualification/resolution function}\\[2pt] Evaluates operation and implementation specifications, actual environment state, requirements and preferences; excludes invalid candidates and ranks eligible ones}}\\[5pt]
$\Downarrow$\\[5pt]
\fbox{\parbox{0.88\textwidth}{\centering\textbf{Selected operation implementation}\\[2pt] For example, CPUGaussianBlur, CUDAGaussianBlur, or MetalGaussianBlur}}\\[5pt]
$\Downarrow$\\[5pt]
\fbox{\parbox{0.88\textwidth}{\centering\textbf{Execution provider and resources}\\[2pt] CPU, Metal, CUDA, OpenCL, Direct3D 12, Vulkan, or a managed runtime}}\\[5pt]
$\Downarrow$\\[5pt]
\fbox{\parbox{0.88\textwidth}{\centering\textbf{New image result or atomic replacement}\\[2pt] Logical semantics are stable; physical allocation is an implementation decision}}
\end{tabularx}
\end{center}

The resolver is one common mechanism, not a custom qualifier and selector implemented by every backend. Its internal eligibility and ranking steps share the same declarative inputs and never allow a preference score to override a failed requirement. The selected operation implementation performs the algorithm; its execution provider supplies the generic execution and resource lifecycle where required.

\subsection{Protocol-library boundary and dependency direction}

The intended package relationship is shown below. The arrow denotes a source-level dependency (imports or consumes protocol contracts), not runtime control flow.

\begin{center}
\begin{tabularx}{0.96\textwidth}{@{}L{0.29\textwidth}Y@{}}
\toprule
\textbf{Library or package} & \textbf{Owns and depends on} \\
\midrule
\texttt{ImageProcessing\allowbreak Protocol} & The primary domain contract: image/layer abstractions, operation definitions/specifications, the generic operation-implementation contract, implementation-specification descriptors, typed composition, and declarative qualification-policy/result contracts. It consumes generic execution facts from \texttt{GeneralExecution\allowbreak Protocol}. \\
\texttt{GeneralExecution\allowbreak Protocol} & The lower-level, reusable execution contract: providers, devices, contexts, generic resources, submission, completion, synchronization, transfer, static capabilities, and observable dynamic execution state. It has no dependency on \texttt{ImageProcessing\allowbreak Protocol}. \\
Concrete execution providers & Implement \texttt{GeneralExecution\allowbreak Protocol} for CPU, Metal, CUDA, OpenCL, Direct3D 12, or another execution technology. \\
Image-operation implementation packages & Supply independent classes conforming to the generic operation-implementation contract. Each class declares an implementation specification and uses a provider or managed algorithm runtime. They are selected and registered by the consuming application. \\
\texttt{ImageProcessingRuntime} & Integrates the two protocols, discovers and snapshots the execution environment, registers heterogeneous implementations, evaluates qualification and routing, plans composed execution, and coordinates result creation or replacement. \\
\bottomrule
\end{tabularx}
\end{center}

This direction avoids two errors. First, a generic execution library must not acquire image-specific concepts merely because image processing is its first consumer. Second, the image-processing API must not expose native GPU mechanisms as its public semantic model. Provider implementations depend on the general protocol; operation adapters depend on the image-processing protocol and whichever provider or library they use. The runtime is the composition point.

\subsection{Decisions recorded in this baseline}

\begin{longtable}{@{}L{0.27\textwidth}Z{0.67\textwidth}@{}}
\toprule
\textbf{Area} & \textbf{Baseline decision} \\
\midrule
Image representation & An image is a logical object/document composed of typed layers; it is not merely a matrix, tensor, or GPU allocation. \\
Layer representation & Raster pixels, masks, categorical labels, and annotations have distinct semantics. They MUST NOT be flattened or numerically interpolated as if they were interchangeable. \\
Storage & Physical pixel storage is separate from logical image/layer identity and is hidden behind a resource abstraction. \\
Operation representation & An operation is an immutable object with stable semantic identity, version, parameters, input/output contracts, and numerical semantics. \\
Operation implementations & One generic implementation protocol is specialized by an associated operation type. Multiple independent classes may implement the same operation. \\
Specifications & Operation specifications describe semantic requirements and permitted behavior; implementation specifications declare the concrete algorithm's support, limits, provider requirements, and resource characteristics. \\
Protocol libraries & \texttt{ImageProcessing\allowbreak Protocol} is the primary domain contract and consumes the lower-level \texttt{GeneralExecution\allowbreak Protocol}; the latter remains generic and image-agnostic. \\
Execution environment & Provider/device capabilities and current availability/state are supplied as environment observations, not embedded in implementation specifications or inferred solely from the operating-system name. \\
Qualification and routing & One common resolver evaluates operation semantics, implementation specifications, environment state, and pipeline policy. It filters invalid candidates and ranks valid candidates. \\
Pipeline policy & Fluent configuration on a function/pipeline expresses hard requirements and soft preferences. Policies are declarative immutable values, not custom per-implementation code. \\
Policy composition & Hard requirements compose conjunctively; incompatible requirements are rejected at builder/composition boundaries with source-aware diagnostics. Preferences use explicit objective dimensions and cannot override requirements. \\
Platform handling & Platform is an environmental fact. No platform-specific preference profile is required to choose Metal on an Apple system or CUDA on a supported NVIDIA system; candidate specifications and discovered availability determine eligibility. \\
Composition & Operations support typed sequential composition. A composed function retains step-scoped policy and supports pipeline-wide objectives. \\
Implementation distribution & Concrete provider and operation-implementation packages are independent of the core protocol/runtime. The application selects and registers them. \\
Result semantics & Producing a new logical image and replacing the state of an existing image are separate explicit API choices. Neither choice guarantees physical in-place memory reuse. \\
Public graph editor & Explicitly excluded from the current design. No node editor, visual graph frontend, or public graph-programming framework is part of this baseline. \\
\bottomrule
\end{longtable}

\section{Problem statement and design requirements}

\subsection{Problem}

The library needs to expose a consistent image-processing model while using different implementation technologies across operating systems and hardware. The same semantic operation may be implemented using Apple's image-processing frameworks, custom Metal kernels, OpenCV CPU code, CUDA, OpenCL, Direct3D 12 compute, or another suitable implementation. Machine-learning inference may be delegated to a separate runtime such as ONNX Runtime or Core ML.

A portable API cannot simply rename the native APIs. These systems differ in resource ownership, scheduling, synchronization, supported formats, precision, and execution granularity. Their commonality exists at the level of operation semantics and execution requirements, not at the level of identical low-level command APIs.

\subsection{Requirements}

The architecture MUST satisfy the following requirements.
\begin{enumerate}
  \item \textbf{Semantic consistency.} An operation has one defined meaning regardless of which implementation executes it.
  \item \textbf{Backend extensibility.} New compute technologies and algorithm libraries can be added without changing the public image and operation semantics.
  \item \textbf{Environment-aware resolution.} A common evaluator uses operation semantics, implementation specifications, provider/device capabilities, current availability/state, and declarative pipeline policy. It must not depend on platform-name heuristics or per-backend qualification code.
  \item \textbf{Multiple implementations.} One operation may have several independent implementation classes; one provider may serve many implementations; and one implementation package may depend on a separate provider package.
  \item \textbf{Layer awareness.} Processing can target an explicitly selected layer or an explicitly rendered composite; unrelated layers are not silently destroyed.
  \item \textbf{Copy/replace semantics.} The caller can request a new image result or replacement of the existing logical image state.
  \item \textbf{Resource efficiency.} The runtime can share immutable resources, use copy-on-write where safe, retain device residency, and avoid unnecessary transfers.
  \item \textbf{Numerical honesty.} Precision, color interpretation, alpha behavior, border rules, and approximation tolerances are part of the operation contract.
  \item \textbf{Testability.} Semantics and qualification rules can be verified independently from real provider implementations and platform-specific hardware, using fake providers and declarative implementation specifications.
  \item \textbf{Failure safety.} A failed replacement MUST NOT leave the source image partially modified.
\end{enumerate}

\subsection{Non-goals}

The design does not attempt to create a universal abstraction over every graphics API feature. It does not require every provider to implement every operation. It does not make every operation differentiable, require automatic differentiation, or convert every input into a generic tensor. It does not make arbitrary Swift closures transparently inspectable. It does not require a node editor, visual graph, or general-purpose workflow engine.

\section{Lessons from existing systems}

The aim is to borrow proven architectural ideas without inheriting entire frameworks or their accidental constraints.

\subsection{OpenCV: image storage and algorithm breadth}

OpenCV's \texttt{Mat} is a useful reference for multidimensional array metadata, row stepping, reference-counted storage, and views. Its copy behavior distinguishes a shallow header copy from an explicit deep clone. This is efficient, but it also demonstrates why object identity, shared storage, and independent storage must be described separately \cite{opencv-mat}.

OpenCV's Graph API (G-API) is an additional architectural reference: it separates public computation definition, graph compilation/optimization, and backend execution \cite{opencv-gapi}. The project may borrow that separation of concerns, but it does not adopt G-API's public graph model as its API foundation.

\textbf{Lesson adopted:} make data layout, stride, ownership, and sharing explicit; reuse mature OpenCV algorithms through adapters where they satisfy the semantic contract.

\textbf{Lesson rejected:} expose \texttt{Mat}, \texttt{UMat}, or \texttt{GpuMat} as the project's public image identity. Those types describe particular storage and execution choices, not the complete semantics of a layered image document.

\subsection{PyTorch: a common numerical object and dispatch}

PyTorch's tensor model describes properties such as shape, data type, layout, strides, and device. Its dispatcher selects kernels through dispatch information and registered implementations; custom backend support is built around registering implementations for dispatch keys \cite{pytorch-tensor,pytorch-dispatch,pytorch-custom-backend}.

\textbf{Lesson adopted:} keep a common semantic operation interface and delegate implementation selection to a runtime dispatch mechanism. Device and layout are relevant facts, but they do not define the domain meaning of image data.

\textbf{Lesson rejected:} make every image a generic tensor. A tensor alone does not state whether values are color channels, fluorescence intensities, a probability map, categorical labels, optical flow, or geometric annotations. The project therefore keeps image/layer semantics above numerical storage and uses tensors where tensor semantics are actually appropriate.

\subsection{ONNX and ONNX Runtime: portable semantics and provider partitioning}

ONNX supplies a versioned representation of operators and typed computational data. ONNX Runtime converts a model into an internal graph, performs provider-independent optimization, partitions it according to execution-provider capability, and assigns supported subgraphs to providers. A default CPU provider can execute work that is not assigned to a specialized provider \cite{onnx-ir,ort-architecture,ort-providers}.

\textbf{Lesson adopted:} provider capability is a qualification question, and implementation selection may occur at subgraph or execution-segment granularity. Correctness and completeness must be preserved when acceleration is unavailable.

\textbf{Lesson rejected:} use ONNX as the public representation for all image processing. It is a model/interchange representation, not the semantic vocabulary for layered image documents and all vision operations.

\subsection{Core Image: deferred image representations}

Apple documents \texttt{CIImage} as a representation or recipe for image processing rather than necessarily materialized pixels. Rendering through \texttt{CIContext} evaluates the processing representation; Core Image can build chains of filters and manage color matching in its working space \cite{coreimage-ciimage,coreimage-cicontext,coreimage-overview}.

\textbf{Lesson adopted:} an operation result need not allocate a distinct pixel buffer after every step. The implementation may retain an internal deferred representation and render or materialize at a suitable boundary.

\textbf{Constraint:} deferred evaluation is an internal optimization unless and until a public lazy-evaluation contract is deliberately introduced. Callers must not receive inconsistent error, lifetime, or completion semantics across providers.

\subsection{Heterogeneous compute APIs are not interchangeable}

Metal exposes resources, command submission, and synchronization mechanisms; NVIDIA CUDA has its own device, stream, memory, and execution model; OpenCL defines contexts, devices, memory objects, command queues, and events; Direct3D 12 assigns explicit responsibilities for command submission, resource state transitions, and synchronization \cite{metal-sync,cuda-guide,opencl-spec,d3d12-queues,d3d12-barriers}. These systems have common concepts but not identical contracts.

\textbf{Lesson adopted:} define a deliberately small common provider contract and optional capability protocols. Do not force every native feature into an artificial lowest-common-denominator API.

\section{Domain model: image, layer, and resource}

\subsection{Image as a logical document}

An \Image{} represents a logical image document. It has a stable identity, a canvas or spatial extent, an ordered collection of layers, and document-level metadata. It is not synonymous with a pixel array. A document may contain original acquisition data, processed rasters, masks, labels, and annotations at the same time.

A conceptual model is:
\begin{lstlisting}
Image
  ImageID
  Revision
  CanvasDescriptor
  OrderedLayers
  DocumentMetadata

Layer
  LayerID
  LayerKind
  LayerProperties
  ContentReference

PixelResource
  ResourceID
  DataDescriptor
  OpaqueStorageHandle
  OwnershipAndLifetime
\end{lstlisting}

The initial Swift domain surface uses immutable value descriptors: \texttt{ImageIdentifier}, \texttt{ImageRevision}, \texttt{ImageCanvasDescriptor}, \texttt{LayerIdentifier}, \texttt{LayerDescriptor}, and \texttt{ImageDocumentSnapshot}. These descriptors establish identity and semantic invariants without prematurely fixing the mutable public \Image{} handle or its actor/transaction implementation.

An \texttt{ImageIdentifier} identifies one logical document. A \texttt{LayerIdentifier} is unique within that document; an external layer reference is the pair \texttt{(ImageIdentifier, LayerIdentifier)}. Two derived documents may retain the same local layer identifier for unchanged logical layers because their image identifiers differ. A \texttt{PixelResource} identity is separate again: two layer references may share one immutable allocation, and one logical layer may refer to a replacement allocation in a later revision.

\texttt{ImageRevision} is monotonically increasing and detects exhaustion rather than wrapping. \texttt{ImageDocumentSnapshot} is an immutable value containing document identity, revision, canvas extent, and ordered layer descriptors. It rejects duplicate layer identifiers and refuses to resolve a layer reference belonging to another image. This snapshot is the input boundary for processing; mutable image ownership, optimistic replacement commit, and concurrency control remain a separate decision recorded in the Architecture Decision Register.

\subsection{Layer kinds and semantics}

A layer's kind determines how its content is interpreted and which operations are valid. The initial semantic vocabulary SHOULD include:
\begin{longtable}{@{}L{0.19\textwidth}L{0.27\textwidth}Z{0.43\textwidth}@{}}
\toprule
\textbf{Layer kind} & \textbf{Representative content} & \textbf{Important constraints} \\
\midrule
Raster & Grayscale, RGB/RGBA, multispectral or floating-point measurements & Has a pixel/data descriptor; color interpretation may be present or absent depending on domain. \\
Mask & Binary or continuous-valued mask & Must distinguish Boolean/binary masks from soft confidence or alpha-like values; thresholding is an explicit operation. \\
Label map & Categorical region or instance identifiers & Values are labels, not intensities. Interpolation that invents new labels is invalid unless a specific label-aware rule is defined. \\
Annotation & Points, paths, contours, boxes, measurements, text anchors & Geometry is expressed in an explicit coordinate space and remains separate from source pixels. \\
\bottomrule
\end{longtable}

This is a semantic classification, not a requirement to create a deep class-inheritance hierarchy. Prefer a small shared layer identity/property contract with typed content representations. Composition and tagged content are often safer than a sprawling subclass tree. A layer type MUST NOT inherit raster-only assumptions merely to reuse a common base class. The initial tagged content model is \texttt{LayerContentDescriptor}: \texttt{raster}, \texttt{mask}, \texttt{labelMap}, or \texttt{annotation}. Mask descriptors require one scientific channel and declare binary, coverage, or probability semantics; a known coverage/probability range must remain within the unit interval. Label-map descriptors require one integer channel explicitly marked as a label identifier and reference a label-set identity. Annotation descriptors identify their schema and coordinate space (pixel, normalized canvas, or a named world coordinate system). These constraints keep label values discrete and annotations separate from raster storage.

Adjustment or procedural layers may be useful later, but they should not be part of the first implementation until evaluation order, parameter persistence, blending, and export semantics are specified.

\subsection{Layer spatial registration}

A raster, mask, or label-map layer has its own sample grid. Matching width/height/depth alone is not enough to prove that two layers refer to the same locations; conversely, different extents can describe the same canvas at different resolutions or through a crop/registration. The image snapshot therefore validates the relation between each pixel layer and the document canvas.

If the layer's \texttt{PixelDataDescriptor.extent} equals the canvas extent and no mapping is supplied, the layer uses the identity mapping. If the extents differ, an explicit \texttt{LayerSpatialMapping} is mandatory. A mapping may also be supplied for equal extents when the layer is translated, rotated, sheared, or otherwise registered differently.

The initial mapping is a validated, invertible 4-by-4 affine transform from layer sample-centre index coordinates to image-canvas sample-centre index coordinates. Coordinates use integer-valued sample centres (the sample at index $(0,0,0)$ is located at that coordinate); the transform is stored row-major and applied to column-coordinate vectors. A two-dimensional layer uses $z=0$. This supports different resolutions, crops, offsets, orientation, scale, and affine registration without assuming the grids coincide.

The mapping describes geometry only. It does not move or interpolate pixel data, does not define a physical/world coordinate system, and does not authorize implicit resampling. An operation that samples one layer on another layer's grid must explicitly choose its interpolation, border, and output-grid semantics. Non-affine deformation fields and calibrated world-coordinate transforms are outside the initial descriptor and require their own explicit model; they must not be approximated by hiding behavior inside this affine mapping.

A snapshot rejects a pixel layer whose extent differs from the canvas when no mapping is supplied. Annotation descriptors retain their own explicit coordinate-space declaration and are not required to have a pixel extent.

\subsection{Layer properties and document composition}

Properties common to renderable layers may include visibility, opacity, compositing mode, coordinate transform, and ordering. Not all properties apply to every content kind. For example, opacity and blend mode can be meaningful for rendered raster content but do not change the categorical identity of a label map. A property that does not apply to a layer kind must be rejected or explicitly defined, not silently ignored.

The image's displayed composite is a rendering result. It is not a replacement for its source layers. Flattening is an explicit destructive or export operation, never an implicit consequence of applying an ordinary operation to a layered image. A composite renderer must use each layer's validated mapping into the document canvas; it must not infer alignment from equal dimensions or silently resample a mismatched layer.

\subsection{Pixel resource and data descriptor}

A \PixelResource{} refers to the numerical content used by a layer. It separates logical content from its physical allocation. A resource may refer to host memory, a Metal texture, a CUDA allocation, an OpenCL memory object, a Direct3D resource, an OpenCV allocation, or an externally owned camera buffer. The semantic descriptor belongs to \texttt{ImageProcessing\allowbreak Protocol}; the opaque execution allocation/handle belongs to \texttt{GeneralExecution\allowbreak Protocol}. The former says what the bytes mean (shape, format, channels, precision, color or scientific interpretation); the latter says how the allocation is owned, accessed, submitted to, synchronized, imported, or released. The image protocol may hold or reference a generic execution resource, but the general execution protocol must never acquire pixel or layer semantics.

The first \texttt{PixelDataDescriptor} is a validated semantic/storage-layout descriptor, not a data allocation. It carries an \texttt{ImageExtent} (width, height, and depth), component type, ordered channel descriptors, data interpretation, and an explicit storage layout. A pixel channel has a stable channel identifier, semantic role, and optional declared numerical value range.

Color and scientific data are separate interpretations. Color interpretation explicitly states color-space identity, channel model, transfer function, and alpha mode. The initial channel-model vocabulary includes grayscale, RGB, YCbCr, and named custom models such as CMYK or Lab; physical channel order is represented separately through the plane mapping. Scientific interpretation may carry a measurement-unit identity and does not invent a color space. The component type (UInt8/UInt16/UInt32, Int16/Int32, FP16/FP32/FP64) does not by itself establish the numerical range or meaning.

Storage layout is either interleaved or planar. Each plane records byte offset, row stride, slice stride, pixel stride, spatial subsampling factors, channel indices, and—when subsampled—the explicit sample origin in layer-grid coordinates. A plane sample at index $(i,j,k)$ maps to layer-grid coordinate $\mathbf{o} + (i s_x,j s_y,k s_z)$, where $\mathbf{o}$ is the declared sample origin and $(s_x,s_y,s_z)$ are the subsampling factors. A subsampled plane without a declared origin is invalid. This records chroma siting explicitly: centered and co-sited chroma are not treated as the same grid merely because both use 2:1 subsampling.

This permits channel order to differ from semantic order (for example BGR), supports planar/subsampled YCbCr with explicit chroma sample locations, and represents volumetric planes. The descriptor verifies that every channel is mapped exactly once, strides can contain each plane, sample origins lie within their subsampling cells, and size arithmetic does not overflow. It calculates the minimum byte capacity for each plane independently and the total capacity required when every plane shares one allocation. A \\texttt{PixelStorageBinding} maps every plane exactly once to a generic \\texttt{ExecutionResource}. Multiple planes may share one contiguous resource, or each plane may use a distinct resource; the binding validates the appropriate per-resource capacity and retains every referenced handle. Plane offsets are relative to the resource assigned to that plane. The binding describes existing storage only: it does not allocate, import, copy, synchronize, or choose a provider.

The descriptor MUST separate at least the following dimensions:
\begin{itemize}
  \item width, height, depth/volume extent, and explicit spatial subsampling per plane;
  \item component type and precision (integer, FP16, FP32, or FP64 where supported);
  \item channel identity, role, channel model, and physical channel mapping;
  \item interleaved or planar storage, byte offset, row/pixel/slice strides, alignment requirements, spatial subsampling, and explicit sample origin/chroma siting;
  \item optional per-channel numerical range and scientific measurement unit;
  \item color-space identity, transfer function, and explicit channel model for color data;
  \item alpha mode (none, straight, or premultiplied) when alpha is present.
\end{itemize}

A \texttt{PixelStorageBinding} is a separate immutable value that associates one or more \texttt{ExecutionResource} handles with the pixel planes. It accepts two common topologies: one resource containing all planes at their declared offsets, or distinct resources for independently stored planes (whose plane-relative offsets are normally zero). A resource group may contain multiple plane indices when those planes share an allocation. Every plane must be assigned exactly once, repeated resource handles must be grouped into one binding record, and the assigned resource capacity must meet the maximum required end offset of its planes. The binding retains resource handles for its own lifetime and exposes generic access declarations for submitted work.

The semantic descriptor does not own the underlying storage. Ownership, mutability, residency, synchronization, and asynchronous lifetime are properties of the resource binding defined through \texttt{GeneralExecution\allowbreak Protocol}. \texttt{PixelDataDescriptor.validateStorageCapacity} verifies the byte capacity of a bound resource against the precomputed minimum; it does not allocate, retain, import, or synchronize that resource.

A color space is not a synonym for a pixel format. An RGBA32F buffer does not, by itself, specify whether RGB is linear or transfer-encoded, whether values are scene-referred or display-referred, or how alpha is interpreted. Scientific intensity data may have no color space at all. The model must allow that without inventing an RGB interpretation.

\subsection{Image versus tensor}

\Image{} and tensor are related but distinct domain types. A tensor is appropriate for general numerical computation and model input/output. An image layer carries additional semantics such as channel meaning, color interpretation, spatial coordinate system, and compositing behavior.

Image-to-tensor conversion MUST be explicit about channel order, layout (for example HWC or NCHW), normalization, precision, scaling, and any color conversion. Model output conversion MUST likewise state whether values are probabilities, logits, labels, coordinates, or another result type. The model runtime must not guess those semantics from shape alone.

\section{Operation objects and semantic contracts}

\subsection{An operation is a declarative object}

An operation is an immutable, parameterized object that names a computation and describes the contract that every implementation must satisfy. It is not a backend method, a kernel, or a function that already knows which implementation will execute.

An operation instance SHOULD expose or reference:
\begin{enumerate}
  \item a stable operation identifier and semantic version;
  \item immutable parameters with validation rules;
  \item typed input and output descriptions;
  \item an \textbf{operation specification} describing semantic requirements and permitted behavior;
  \item data constraints (formats, dimensions, layouts, layer kinds, and parameter domains);
  \item numerical semantics (input/storage precision, computation precision, accumulation precision, output precision, rounding, range, color/alpha rules, border/interpolation rules);
  \item mutation, aliasing, and effect semantics.
\end{enumerate}

For example, a Gaussian blur operation is not fully specified by \textquotedblleft{}sigma = 1.2\textquotedblright{}. Its border extension, kernel construction/truncation, channel treatment, alpha treatment, and numerical contract are also relevant. Two implementations carrying the same operation identity but disagreeing on these choices are not conformant merely because both are called Gaussian blur.

\subsection{Operation specification versus implementation specification}

These specifications are complementary and have different owners.

\textbf{Operation specification} describes the operation's semantic truth: what result is required, which input domains and parameters are valid, what numerical guarantees apply, and which implementation variations are permitted. It may be refined by the concrete operation parameters and by the caller's pipeline requirements.

\textbf{Implementation specification} describes one concrete implementation class: which operation identity/version it implements, what semantic modes and parameter ranges it supports, which input/output representations and precision modes it can realize, which provider capabilities it requires, and its resource/execution characteristics. It is declarative metadata, not a custom \texttt{qualify()} function.

For a \texttt{GaussianBlur}, an implementation specification may declare supported border modes, supported component formats, computation/accumulation precision, maximum radius, required device capabilities, and resource limits. \texttt{CUDAGaussianBlur} declares a CUDA provider requirement; \texttt{MetalGaussianBlur} declares a Metal requirement; \texttt{CPUGaussianBlur} declares the CPU capability requirements. The runtime compares these declarations with the actual environment. It does not encode special-case logic such as \textquotedblleft{}if macOS then prefer Metal\textquotedblright{}.

Implementation specifications MUST distinguish:
\begin{itemize}
  \item static, tested support declared by the implementation;
  \item provider/device capabilities reported by the execution provider;
  \item transient availability and load, reported by the current environment snapshot.
\end{itemize}
A specification must not claim that a device is currently available or currently idle. Those are runtime facts.

\subsection{One generic operation-implementation contract}

All concrete algorithm classes conform to one generic implementation contract specialized by an associated operation type. Conceptually:

\begin{lstlisting}
public protocol ImageOperationImplementation {
    associatedtype Operation: ImageOperation

    var specification: ImplementationSpecification { get }

    func execute(
        operation: Operation,
        input: Operation.Input,
        context: ExecutionContext
    ) async throws -> Operation.Output
}
\end{lstlisting}

The protocol has no mandatory per-class qualification method. The implementation publishes its specification and executes only after the runtime has resolved it as a valid candidate. A future operation-specific extension may be introduced only when a real semantic constraint cannot be represented declaratively; it must be owned by the operation/runtime contract, not duplicated as arbitrary qualification code in every backend class.

The associated operation type preserves type safety for the public implementation API. The runtime may use private type erasure to register different operation types in one heterogeneous registry; that mechanism must not leak as \texttt{Any}-based application code.

\subsection{Operation identity and versioning}

An operation identifier names its semantics, not the implementation library. An operation called \texttt{image.gaussian-blur} remains that operation whether implemented by a custom Metal kernel, OpenCV, CUDA, OpenCL, or another library.

The semantic version changes when a change can alter observable results or accepted inputs. Implementation version and provider version are separate metadata. This separation supports reproducibility, diagnostics, and future serialization of parameterized operations without binding stored processing descriptions to one vendor library.

Operation identity and versioning SHOULD be established before persistent operation serialization is implemented. A serialization format is not a prerequisite for the first runtime, but stable identifiers avoid painting the API into a corner.

\subsection{Typed inputs and outputs}

The operation model must not assume that every operation is a unary image-to-image transformation. Examples include:
\begin{align*}
  &\operatorname{GaussianBlur}: \texttt{RasterLayer} \rightarrow \texttt{RasterLayer},\\
  &\operatorname{Resize}: \texttt{RasterLayer} \rightarrow \texttt{RasterLayer},\\
  &\operatorname{Blend}: (\texttt{ImageContent},\texttt{ImageContent}) \rightarrow \texttt{RasterLayer},\\
  &\operatorname{Segment}: \texttt{Image} \rightarrow \texttt{SegmentationResult},\\
  &\operatorname{Measure}: \texttt{Image} \rightarrow \texttt{MeasurementResult}.
\end{align*}

For composition, a multi-input operation can receive a typed input record (for example, foreground, background, mask, and blend parameters). This preserves type clarity without requiring a public general graph abstraction. A segmentation or measurement result may subsequently be attached to an image as a corresponding layer through an explicit image-editing operation.

\subsection{Pure transformations and explicit effects}

Ordinary image transformations SHOULD be pure with respect to their input resources: given the same immutable input data, operation version, and parameters, they produce the same defined result within the operation's numerical contract. Implementations MUST NOT mutate source resources behind the caller's back.

Not every useful operation is pure. Camera acquisition, file I/O, video tracking across frames, random sampling, and operations that maintain temporal state have effects or state. They must be represented as effectful or stateful operations with explicit inputs, outputs, lifecycle, and failure semantics. They must not be mislabeled as ordinary pure image transforms.

The project does not need a general effect-system framework at the outset. A small effect classification in the operation contract is sufficient until concrete use cases justify more.

\section{Composition and image result semantics}

\subsection{Typed composition}

If operation $A$ maps type $X$ to type $Y$, and operation $B$ maps $Y$ to $Z$, then their sequential composition maps $X$ to $Z$:
\[
  A : X \rightarrow Y,\qquad B : Y \rightarrow Z
  \quad\Longrightarrow\quad
  B\circ A : X \rightarrow Z.
\]

The Swift-facing API may express this using a method such as \texttt{then}: the first operation's output contract must satisfy the second operation's input contract. The composition object records the ordered operations and their parameter instances. It is immutable once created or executed.

Illustrative API sketch (not a frozen signature):
\begin{lstlisting}
let correction = ColorCorrection(exposure: 0.25, contrast: 1.1)
let blur        = GaussianBlur(sigma: 1.2, border: .reflect)
let resize      = Resize(width: 1024, height: 768)

let pipeline = correction.then(blur).then(resize)
    .qualifying {
        .require(.precision(.computation, .fp32))
        .prefer(.lowLatency)
    }

let result = try await source.applying(
    pipeline,
    to: .layer(sourceLayer.id)
)
\end{lstlisting}

The public abstraction is a composed operation, not a visual graph. The runtime may create an internal execution plan to optimize the composition, but that internal plan is not a node-editor API and does not become an additional public programming model.

\subsection{Composition boundary and multi-input operations}

Sequential composition is the first required composition mechanism. Multi-input operations such as blending are represented through explicit typed input values. A general-purpose branching graph, arbitrary closures, implicit fan-out, or graph rewriting is not required for the first version.

If later requirements prove that users need reusable branches or shared intermediate values, those can be evaluated against concrete workloads. They should not be introduced merely because other frameworks provide graph APIs.

\subsection{New image versus replacing the input image}

The public API MUST distinguish a new logical result from replacement of the current logical image state. These are semantic choices, not allocation hints.

Conceptual API:
\begin{lstlisting}
// Produce a new logical image. Source remains unchanged.
let result = try await image.applying(operation, to: target)

// Update the existing logical image.
try await image.apply(operation, to: target)
\end{lstlisting}

The exact names may change, but the two behaviors must remain distinguishable. A Boolean such as \texttt{copy: true/false} is discouraged because it obscures whether the flag controls document identity, object mutation, or physical storage reuse.

\subsubsection{New logical image}

A new-image operation creates a new logical image identity and a result state. Unchanged layers and immutable resources MAY be shared. The resource system can use copy-on-write when a later operation requires writable storage. The new image is logically independent even when some physical resources are shared.

The identity policy for copied layer identifiers must be explicit. A useful model is to scope a layer reference by its owning image identity, so that an unchanged layer identifier can be reused within a new document without creating an ambiguous global reference. This detail must be validated before IDs become part of a persisted file format.

\subsubsection{Replacement}

Replacement preserves the logical identity of the existing image while committing a new revision of its state. It does not promise reuse of the original allocation. A GPU kernel may require a new output texture; a resize necessarily changes dimensions; a framework may only produce an immutable result. The implementation can allocate new storage and then replace the logical resource reference.

Replacement MUST be transactional:
\begin{enumerate}
  \item resolve and qualify a valid implementation;
  \item execute into separate or otherwise safely managed output resources;
  \item validate output descriptors and operation postconditions;
  \item atomically publish the new image revision;
  \item release obsolete resources only after all readers and device work have completed.
\end{enumerate}

If execution or validation fails before commit, the original logical image state remains intact. A direct physical in-place overwrite is permitted only when aliasing, failure atomicity, and synchronization rules prove that it is safe; it is an optimization, not the meaning of replacement.

\subsection{Operation target: layer or composite}

An operation that acts on image content MUST make its target semantics explicit. At minimum, distinguish:
\begin{itemize}
  \item a specific layer (for example, blur this raster layer);
  \item a rendered composite (for example, blur the visible composite of selected layers);
  \item a set of explicitly supplied inputs (for example, blend foreground and background layers).
\end{itemize}

A layer-targeted operation preserves unrelated layers. A composite-targeted operation processes a rendering result and produces a new raster result; it must not silently flatten and discard the source layer structure. Composite rendering rules, color-management space, opacity, and blend behavior must be explicit.

\section{Protocol libraries and execution-provider architecture}

\subsection{Protocol-library dependency direction}

The architecture is not organized by placing a generic compute abstraction above the image-processing API. The domain-facing \texttt{ImageProcessing\allowbreak Protocol} library is the primary contract. It defines what an image-processing operation is, what an implementation must promise, and how a candidate is qualified for a particular image and operation. It depends on the lower-level \texttt{GeneralExecution\allowbreak Protocol} only for execution facts and contracts that are genuinely generic.

\texttt{GeneralExecution\allowbreak Protocol} is deliberately domain-neutral. It defines how a provider exposes devices, contexts, resources, work submission, completion, synchronization, and transfers. It MUST NOT define images, layers, pixel meaning, blur, blending, segmentation, or image-specific qualification. That keeps it reusable by other domains and prevents an image-processing library from becoming the owner of generic compute infrastructure.

The direction is therefore:
\[
\begin{gathered}
  \texttt{ImageProcessing\allowbreak Protocol}\\
  \downarrow\;\text{depends on}\\
  \texttt{GeneralExecution\allowbreak Protocol}\\
  \downarrow\;\text{implemented by}\\
  \text{concrete execution-provider implementations}
\end{gathered}
\]
The arrow is a dependency direction. Concrete provider packages implement the lower-level protocol. Image-processing adapters implement the domain operation contracts and use provider implementations to execute them. The runtime integrates the two and is not itself a hardware provider.

\subsection{GeneralExecutionProtocol library}

The general execution library exposes only concepts that remain meaningful outside image processing. A provider supplies access to a compute substrate, provider/device discovery, generic resources, and optional work execution capabilities. Examples include CPU execution, Metal, CUDA, OpenCL, Direct3D 12, and potentially Vulkan.

The common provider contract SHOULD cover:
\begin{itemize}
  \item provider identity, version, discovery, and provider-level capabilities/state;
  \item device identity, discovery, device-level capabilities/state, and execution-context creation;
  \item a stable execution-target identity, including provider-managed targets whose internal devices are opaque;
  \item typed generic resource identity, descriptors, allocation requests, ownership/lifetime, import/export, and residency;
  \item asynchronous work dependencies, typed completion, cancellation, and synchronization where supported;
  \item structured errors and observability for provider/device failures.
\end{itemize}

\textbf{Generic resources.} An \texttt{ExecutionResource} is an opaque, reference-counted provider allocation. Its generic descriptor identifies its owning execution target, resource kind (for example buffer, texture, or provider-defined opaque kind), capacity, alignment, and memory domain. It does not describe image dimensions, pixel formats, channel order, color spaces, tensor layout, or scientific units. Those belong to the domain consuming the resource. Resource identifiers are scoped by execution target; the pair \texttt{(target, identifier)} is the stable generic resource reference.

Allocation and import are optional context capabilities. An allocation request states minimum capacity, alignment, kind, and memory domain; a provider must verify that its returned descriptor satisfies that request and belongs to the context's target. Import produces a context-local resource handle. The imported handle MUST retain the source handle for as long as the import or any submitted work uses it. This makes the lifetime relationship explicit without exposing native allocation pointers or introducing a second manual-release system on top of reference ownership.

\textbf{Work submission.} Providers MAY expose an \texttt{ExecutionWorkSubmittingContext}. A typed \texttt{ExecutionWork} declares resource accesses (read, write, or read/write) and completion dependencies, and executes using the context. Submission returns a typed \texttt{ExecutionSubmission} whose value/completion operation is the public boundary at which the result is guaranteed valid. The provider MUST retain the work item and all declared resources until completion, even if the caller releases its local references. Dependencies MUST complete successfully before dependent work starts.

Cancellation is a request, not a universal promise that already-submitted hardware work can be interrupted. The result distinguishes cancellation before execution, cancellation requested after execution may have begun, and work that is already too late to cancel. After requesting cancellation, the caller still observes completion before considering the work's resources reusable. Provider implementations map these rules onto their native mechanisms.

\textbf{Synchronization.} Explicit context-wide or resource-scoped synchronization is an optional capability. Waiting for one submission is preferred when that is sufficient. Providers retain their native event, fence, barrier, queue, and command-buffer mechanisms; the common contract does not define a universal command object.

These contracts are intentionally capability-based. A provider may support allocation and import but not explicit synchronization, or support typed work submission while keeping queue management private. An operation implementation declares the capabilities it requires; the runtime must not assume every provider implements every optional interface.

\textbf{Context capabilities.} Each \texttt{ExecutionContext} exposes its own generic \texttt{ExecutionCapabilities} in addition to the provider/device descriptors. Standard capability identifiers cover buffer/texture resources, allocation/import, work submission/dependencies, synchronization, and peer-resource transfer. The \texttt{ExecutionCapabilityScope} can target provider, device, provider-or-device, or the created context. Context capability reports are the exact execution target's facts; they may be narrower than the device's static capabilities.

A context MUST report a standard capability only when it actually supports the corresponding optional context protocol and operation. The optional protocol conformance is the executable proof; the capability descriptor is the declarative fact used by implementation qualification. A context that does not expose capability reporting receives an empty descriptor by default, so it cannot accidentally satisfy a required context capability. Standard identifiers are namespaced under the generic execution capability vocabulary, so implementation packages do not invent incompatible strings for common requirements. The descriptor may also include provider-defined capability identifiers for features not standardized by \texttt{GeneralExecution\allowbreak Protocol}.

The resolver first evaluates provider/device requirements against the environment snapshot. When a concrete candidate context is prepared, it evaluates any context-scoped requirements against that context's capability descriptor and verifies the required optional protocol conformance before submitting work. A context-level mismatch rejects that binding before the operation executes; the runtime may then consider another eligible candidate within the original hard constraints.

The common contract MUST NOT force every provider to expose a Metal-like command buffer, a CUDA-like stream, an OpenCL command queue, or a Direct3D command list as the same public object. The common API expresses resource and completion intent; provider-specific implementations retain native details. Optional capability protocols may cover specialized buffer/texture allocation, external-memory interop, peer-device transfer, native kernel encoding, asynchronous queues, and specialized synchronization.

\subsection{ImageProcessingProtocol operation and implementation contracts}

The image-processing protocol library defines the domain-facing operation and implementation contracts. An operation object describes meaning and parameters. Its specification describes required and permitted semantics. A concrete implementation class conforms to the generic \texttt{ImageOperationImplementation} contract with one associated operation type and publishes an \texttt{ImplementationSpecification}.

This is a many-to-many relationship at the system level:
\begin{itemize}
  \item one operation type may have several implementation classes;
  \item one execution provider may support many operation implementation classes;
  \item an operation implementation may use a provider, a higher-level algorithm library, or a managed execution runtime;
  \item a consuming application may register only the implementation packages it chooses to deploy.
\end{itemize}

Each implementation specification SHOULD declare:
\begin{enumerate}
  \item operation identifier and supported semantic version range;
  \item implementation identifier, version, and provenance;
  \item execution binding: required provider/runtime and required generic capabilities;
  \item supported input/output layer kinds, formats, layouts, dimensions, and parameter ranges;
  \item supported input/storage, computation, accumulation, and output precision modes;
  \item supported semantic options such as border modes, interpolation, color spaces, alpha handling, and deterministic behavior;
  \item resource requirements, access modes, memory constraints, and import/export needs;
  \item static execution characteristics or cost-model hints, clearly distinguished from measured current state.
\end{enumerate}

An implementation must not redefine the operation's meaning. It may declare that it supports only a subset of the operation's semantic domain. A backend does not need to implement every operation; an operation does not need an implementation on every provider. No backend package is required to modify the operation type or runtime when adding a new conformant implementation.

\subsection{The common qualification/resolution contract}

The image-processing domain owns the semantics and policy vocabulary for qualification. The runtime supplies the available implementations and a current generic execution-environment snapshot. A single common resolver evaluates the four complementary inputs: operation and its specification, implementation specification, execution environment, and pipeline qualification policy.

The common evaluator is data-driven. It compares declarative specification fields and generic provider capabilities; it does not contain a growing list of special cases for CPU, CUDA, Metal, OpenCV, or particular operating systems. The resolver is described in detail in the qualification section.

\subsection{Native providers and library adapters}

Not every library is a low-level compute provider. The architecture must distinguish hardware APIs from algorithm libraries and model runtimes.

\begin{longtable}{@{}L{0.21\textwidth}L{0.27\textwidth}Z{0.42\textwidth}@{}}
\toprule
\textbf{Technology} & \textbf{Architectural role} & \textbf{Integration interpretation} \\
\midrule
CPU & Execution provider & Baseline portable implementation and fallback where supported. \\
Metal & Compute provider & Native Apple GPU execution, resources, command submission, and synchronization. \\
CUDA & Compute provider & NVIDIA GPU execution and resource/stream management. \\
OpenCL & Compute provider & Heterogeneous device execution through contexts, memory objects, command queues, and events. \\
Direct3D 12 & Compute provider & Windows GPU execution through command queues/lists, resources, barriers, and fences. \\
Vulkan & Potential compute provider & Add only when a concrete supported implementation and test matrix justify it. \\
Core Image & Higher-level algorithm/runtime adapter & Provides image operations and manages its own evaluation context. Treat it as a registered implementation/runtime; do not assume all resources are directly interchangeable with raw Metal resources. \\
Accelerate/vImage & Algorithm implementation adapter & CPU/vectorized image operations. It may execute without a separately visible GPU provider. \\
OpenCV & Algorithm library adapter & Register concrete CPU, CUDA, or OpenCL implementations separately where the actual build and data path support them. \\
ONNX Runtime & Model execution runtime & Treat a model/session or supported model subgraph as an implementation boundary; its own execution providers partition internal model operations. \\
PyTorch & Tensor/model runtime & Use when needed for tensor computation or model execution; do not make it a dependency of the core image API. \\
Core ML & Model execution runtime & A model implementation adapter with platform-specific device selection and constraints. \\
\bottomrule
\end{longtable}

An adapter may manage its own compute context and scheduling. In that case, the outer runtime must treat the adapter as an opaque execution candidate unless the adapter exposes reliable resource and partition information. It must not double-schedule internal operations that the external runtime already owns.

For example, an ONNX Runtime session may partition a model internally among execution providers. The image runtime should normally qualify and select the model/session as one higher-level implementation candidate, rather than attempting to select a backend independently for every internal ONNX node. The same rule applies to a Core Image filter chain when Core Image owns evaluation.

\subsection{General execution capability and state model}

A provider capability descriptor reports facts, not promises inferred from a platform name. Depending on the provider, useful static facts include:
\begin{itemize}
  \item provider/device identity, vendor, architecture, driver/runtime version, and supported platform/runtime combinations;
  \item supported generic resource kinds, opaque native resource-format identifiers, import/export paths, and alignment constraints;
  \item standardized generic execution capabilities reported by providers and devices;
  \item context-specific capabilities reported only by the actual prepared execution context;
  \item supported compute features and precision modes;
  \item maximum dimensions, allocation limits, and device memory characteristics;
  \item asynchronous submission, synchronization, and interoperability capabilities;
  \item whether scheduling and execution are externally visible or managed internally by a higher-level runtime.
\end{itemize}

The environment also exposes dynamic observations where the provider can report them:
\begin{itemize}
  \item device health and current availability;
  \item available memory and memory pressure;
  \item queue depth, outstanding work, and estimated queue delay;
  \item workload/utilization observations where the platform exposes trustworthy metrics;
  \item thermal/power state or throttling information where available;
  \item current resource residency and outstanding synchronization dependencies.
\end{itemize}

Dynamic observations MUST carry observation time and validity/freshness information. The protocol must not pretend that every operating system or driver exposes a reliable global GPU-utilization number. When such telemetry is unavailable, the runtime can still use its own outstanding submissions, queue depth, allocation knowledge, and measured execution history. Missing telemetry is represented as unknown, not as zero load.

The operating-system platform is an environmental fact and may constrain which provider packages can load or which capabilities exist. It is not itself a preference rule. A CUDA implementation is eligible only when its declared provider/capability requirements are satisfied by an actually discovered CUDA execution path; a Metal implementation is eligible only when a compatible Metal path exists. The resolver need not encode a platform-to-provider preference table.

\section{Qualification and implementation resolution}

\subsection{One common resolver; distinct internal stages}

The caller configures a function/pipeline declaratively. The runtime invokes one common qualifying/resolution function over the composed operations and the registered implementation specifications. The resolver is smart about implementation declarations and current hardware state; individual implementation classes do not each implement their own qualification algorithm.

Conceptually:
\[
  R(o, x, s_o, s_i, e, k, p) \longrightarrow
  \text{ranked candidate evaluations and an execution choice},
\]
where:
\begin{itemize}
  \item $o$ is the operation instance and semantic version;
  \item $x$ is the concrete input, target, and resource descriptor/state;
  \item $s_o$ is the operation specification and its semantic requirements;
  \item $s_i$ is a candidate implementation specification;
  \item $e$ is the execution-environment snapshot, including provider/device capabilities and current state observations;
  \item $k$ is the candidate-specific execution-context capability descriptor, obtained when the candidate context is prepared;
  \item $p$ is the immutable qualification policy attached to the function/pipeline.
\end{itemize}

The public architecture has one resolution entry point. Internally, it must preserve two logically distinct stages:
\begin{enumerate}
  \item \textbf{Static eligibility:} reject candidates that fail operation semantics, implementation support, input/output, precision, provider/device, or hard-policy requirements using the environment snapshot.
  \item \textbf{Context validation:} prepare a candidate context for remaining bindings, then check context-scoped capability requirements and the corresponding optional protocol conformance before any operation work is submitted.
  \item \textbf{Ranking:} compare only fully eligible candidates using preferences, current state, and estimated complete execution cost.
\end{enumerate}
This distinction is not an invitation for each backend to implement its own qualifier and selector. It is an internal structure of one shared, data-driven resolver. A score can never make an ineligible candidate eligible.

\subsection{The four complementary inputs}

\begin{longtable}{@{}L{0.20\textwidth}L{0.30\textwidth}Z{0.39\textwidth}@{}}
\toprule
\textbf{Input} & \textbf{Owner} & \textbf{What it contributes} \\
\midrule
Operation specification & \texttt{ImageProcessing\allowbreak Protocol} & Semantic meaning, valid parameters/inputs, required guarantees, and permitted implementation variation. \\
Implementation specification & Concrete operation implementation package & Static, tested support for the operation, provider requirements, precision behavior, modes, limits, and resource characteristics. \\
Execution environment & \texttt{GeneralExecution\allowbreak Protocol} providers plus runtime observations & Actual platform, loaded providers, discovered devices, provider/device capabilities, resource availability, residency, queue pressure, state freshness, and prepared candidate-context capabilities. \\
Pipeline policy & Immutable policy attached through the functional API & Caller requirements and optimization preferences, globally or scoped to an operation/step. \\
\bottomrule
\end{longtable}

These inputs are complementary, not competing preference sets. The operation specification says what must be true. The implementation specification says what one algorithm can provide. The environment says what can be used now. The pipeline policy says which valid outcomes the caller requires or prefers.

\subsection{Declarative pipeline qualification policy}

The functional pipeline MAY expose a fluent configuration extension that builds an immutable policy value. It configures the common resolver; it does not inject code into an operation or implementation class.

Illustrative API (names and overloads are provisional):

\begin{lstlisting}
let pipeline = ImagePipeline {
    GaussianBlur(sigma: 1.5)
    ColorCorrection(...)
    Resize(width: 1024, height: 768)
}
.qualifying {
    // Hard requirement applied to applicable numeric operations.
    .require(.precision(.computation, .fp32))

    // Soft objective: rank valid candidates for end-to-end latency.
    .prefer(.lowLatency)

    // Optional rule scoped to a particular operation type/step.
    .for(GaussianBlur.self) {
        .require(.borderMode(.reflect))
    }
}
\end{lstlisting}

A pipeline-wide requirement is inherited by all applicable operations. An operation-scoped policy narrows or refines the policy for that operation occurrence/type. Composition preserves step-scoped policy and does not leak one operation's local preference into unrelated steps. The builder validates that each policy dimension applies to its target operations; ambiguous applicability must produce a diagnostic rather than be silently ignored.

The policy is independent of concrete provider names by default. The caller normally states a requirement or objective such as FP32 computation, low latency, low memory use, deterministic results, or bounded energy. The resolver discovers which implementation/provider combination best satisfies it. An explicit provider restriction MAY exist for reproducibility or deployment control, but it is a hard constraint, not a platform-conditioned preference.

\subsection{Requirements and preferences are different}

A \textbf{requirement} removes candidates that cannot satisfy it. A \textbf{preference} ranks candidates that remain.

For example:
\begin{itemize}
  \item \texttt{require(computationPrecision: fp32)} means an FP16-only implementation is not eligible.
  \item \texttt{prefer(computationPrecision: fp32)} means FP32 is ranked higher, but a lower precision may be selected only if the operation specification permits it and its declared numerical guarantee remains acceptable.
  \item \texttt{prefer(lowLatency)} may select CPU while a GPU is busy if CPU is eligible and its estimated end-to-end latency is lower.
  \item \texttt{require(provider: CUDA)} forbids CPU fallback. If CUDA is unavailable or cannot reserve required resources, the request waits/retries according to the execution policy or returns a structured failure.
\end{itemize}

Precision MUST be dimensioned. Input/storage precision, computation precision, accumulation precision, and output precision are different properties. A bare \texttt{FP32} label is ambiguous. The policy API should therefore identify the precision dimension it constrains. A preference never authorizes a result outside the operation's numerical contract; a hard requirement is needed when the caller specifically requires that precision.

\subsection{Environment-aware qualification and routing}

The resolver must consider more than whether a GPU exists. It evaluates:
\begin{enumerate}
  \item provider and device compatibility with the implementation specification;
  \item operation-specific format, parameter, border, color, alpha, and precision support;
  \item memory capacity, current memory pressure, resource importability, and residency;
  \item device health, availability, queue depth, outstanding work, and estimated delay where observable;
  \item conversion, allocation, upload/download, synchronization, and output-materialization costs;
  \item pipeline hard requirements and soft objectives.
\end{enumerate}

A busy GPU is not automatically semantically unqualified. It may remain an eligible candidate with a worse cost estimate. If the pipeline prioritizes latency and an eligible CPU route is predicted to finish sooner, the resolver can select CPU. If the pipeline requires CUDA, the CPU route is ineligible regardless of its speed. If the GPU is unavailable, lost, or cannot satisfy a mandatory resource requirement, the candidate is temporarily unavailable or rejected with a structured reason.

\textbf{Context-level requirements.} Provider/device capabilities are not sufficient when an operation implementation requires a context interface such as resource allocation, work submission, dependency handling, or explicit synchronization. The runtime prepares a candidate context only after static provider/device checks pass, then verifies the context-scoped capability constraints against \texttt{context.capabilities} and verifies that the corresponding optional context protocols are actually implemented. This is part of the common resolver's eligibility process, not a backend-specific qualifier. If context creation or capability verification fails, the candidate is rejected with a structured reason before any operation work is submitted.

The runtime must not infer that a GPU is idle because no global utilization metric is available. It uses the best reliable observations: its own queue/outstanding-work accounting, provider-reported state, available memory, measured latency, and telemetry freshness. Every dynamic observation SHOULD include an observation timestamp and, where practical, a confidence or source quality.

\subsection{Conflict prevention and policy composition}

The fluent builder produces a normalized, immutable policy rather than an unstructured list of positive and negative predicates. Each constraint dimension has a well-defined representation: for example, an allowed precision set, a required provider, or one explicit optimization objective/profile. This prevents contradictory declarations from accumulating unnoticed.

Hard constraints compose conjunctively. If two composed functions permit sets $A$ and $B$ for the same dimension, the combined admissible set is:
\[
  A_{\mathrm{combined}} = A \cap B.
\]
An empty result is a policy conflict and MUST be rejected when the builder is finalized or when the functions are composed. For example, requiring FP32 in one constituent and FP16-only in another, when both constraints apply to the same operation scope, is invalid. A requirement and an exclusion of the same provider is also invalid.

Preferences are not intersected as sets. They are represented as explicit objective dimensions and, when multiple objectives must be traded off, by a declared priority or weight model. Contradictory declarations for the same objective at the same scope are rejected unless the caller explicitly supplies a valid ordering/weighting. Composition preserves preferences scoped to individual steps; pipeline-wide objectives apply to the complete execution plan. An operation-local rule cannot silently erase a pipeline-wide hard requirement.

Conflict diagnostics MUST preserve provenance: policy scope, pipeline/operation/step identity, source declaration, conflicting declaration, and the incompatible dimension. Validation occurs:
\begin{itemize}
  \item while building/finalizing a policy;
  \item when composing functions or pipelines;
  \item when loading a serialized policy;
  \item again at runtime as a defensive boundary for dynamic or externally supplied declarations.
\end{itemize}
Runtime availability failures are not policy conflicts. Requiring CUDA on a machine without an available compatible CUDA device is an unsatisfied requirement, not a contradictory configuration.

\subsection{Cost model and candidate ranking}

For each eligible candidate, the resolver estimates complete cost rather than comparing kernel duration alone. A useful initial model is:
\[
  C_{\mathrm{total}} =
  C_{\mathrm{queue}} +
  C_{\mathrm{transfer}} +
  C_{\mathrm{conversion}} +
  C_{\mathrm{allocation}} +
  C_{\mathrm{context}} +
  C_{\mathrm{compute}} +
  C_{\mathrm{synchronization}} +
  C_{\mathrm{materialization}}.
\]
The model may also account for peak memory, power/thermal objectives, expected throughput, and current resource pressure.

Consider a case in which the GPU kernel is much faster but the GPU queue is busy and the input is resident in host memory. CPU can be the better candidate because it avoids queue delay and upload. Conversely, a GPU may be preferable when input is already resident there and a sequence can remain on the same device. This is why routing must consider resource residency and the composed sequence.

Initial estimates may be simple and partly measured. The architecture requires honest cost boundaries and observable decisions, not an elaborate predictive model before evidence exists. Unknown telemetry must remain unknown and must not be silently treated as zero cost.

For a composed function, routing must consider compatible implementation sequences or execution segments, not choose each operation independently by its isolated kernel estimate. The planner accounts for the cost of transitions between candidates, resource residency, and intermediate materialization. A locally fastest implementation can be globally slower when it forces transfers or synchronization before the next operation. The initial planner may use a simple segment-based strategy; it does not need a general optimizer or a public graph representation.

\subsection{Resolution, reservations, and changing device state}

A qualification result is based on an environment snapshot and is not a permanent reservation. Device load, available memory, and resource availability can change between resolution and execution.

Before execution, the coordinator MUST acquire/reserve the resources required by the selected plan and validate critical availability assumptions. If reservation fails or the snapshot is stale, the runtime may refresh the environment and resolve again, but only within the original hard constraints. It must not silently fall back to a forbidden provider or lower precision. If fallback is permitted, it may resolve to another eligible candidate. If no candidate remains, return a structured availability/resource failure.

The resolver should return the ranked candidate evaluations, selected candidate/plan, rejected/temporarily-unavailable reasons, cost assumptions, and environment-snapshot identity/time. This supports diagnostics and reproducibility.

A candidate is an executable binding, not merely an implementation class: it consists of an implementation instance, the provider or managed-runtime instance it uses, the selected device/execution target, and the environment snapshot against which it was evaluated. One generic implementation class can therefore produce multiple candidates when several compatible devices are discovered. A managed runtime whose internal device partitioning is opaque is represented as one candidate unless it exposes reliable partition/resource details.

\subsection{Fallback policy}

CPU fallback is not a universal rule. It is one possible route selected when:
\begin{itemize}
  \item a CPU implementation is registered and qualified for the same operation semantics;
  \item the pipeline does not require a different provider;
  \item the CPU route satisfies precision, output, and resource requirements;
  \item its estimated cost is competitive for the selected objective.
\end{itemize}

No silent downgrade of precision, color space, alpha semantics, border behavior, label handling, or deterministic guarantees is permitted. A preference can be compromised; a requirement cannot.

\section{Execution, memory, and lifecycle}

\subsection{Logical resources and physical resources}

A \PixelResource{} is a logical handle whose physical representation belongs to a provider or an external owner. The handle should expose a stable descriptor, lifetime and access contract while keeping provider-native pointers and objects behind an adapter boundary.

A resource may be:
\begin{itemize}
  \item host-resident;
  \item resident on one GPU/device;
  \item shared through a verified interop mechanism;
  \item externally owned and temporarily leased;
  \item represented lazily by a higher-level image runtime until materialization.
\end{itemize}

The architecture MUST NOT assume that unified memory means identical access cost, ordering, or lifetime across APIs. Nor may it promise zero-copy for every provider pair. A transfer-free path exists only where both providers expose a valid interop mechanism with compatible formats and synchronization.

\subsection{Ownership and lifetime}

The generic execution resource handle is a reference-counted opaque allocation. The provider implementation owns the native storage; callers own references to the handle. A provider may release or recycle the native allocation only after the resource handle and all outstanding submission/import references are gone.

An imported resource is a new handle bound to the importing execution target. It MUST retain the source resource for as long as the imported handle or submitted work refers to it. This applies to externally owned buffers as well: the source handle represents the external owner's lease/retention rule, and the imported handle must not outlive that lease silently.

A submitted work item declares each resource and its access mode. The provider MUST retain the work item and every declared resource until the submission reaches completion, even if the caller drops its references or requests cancellation. A cancellation request does not authorize early release. The caller observes completion (success, failure, or cancellation) before considering scratch storage or external buffers reusable.

These rules allow reference ownership to express the normal lifetime path. A separate manual release call is not required by the common contract; provider wrappers may still expose explicit pool/lease operations when the native API requires them, but they must preserve the same lifetime invariant.

\subsection{Read, write, and alias contracts}

Every operation implementation must declare how it accesses its inputs and outputs. The initial contract should distinguish:
\begin{itemize}
  \item read-only input;
  \item newly allocated output;
  \item output permitted to alias an input under defined conditions;
  \item exclusive writable input where mutation has been explicitly requested;
  \item temporary scratch resources.
\end{itemize}

The default for ordinary image transformations is read-only input and a logically independent output. Physical aliasing is an optimization that requires proof that no observable source mutation can occur and that concurrent readers remain safe.

\subsection{Synchronization and concurrency}

The common execution model represents resource accesses, explicit completion dependencies, and optional resource/context synchronization. A submission can depend on prior completion handles; a provider must not begin dependent work until all dependencies succeed. A failed or cancelled dependency propagates as a failure of the dependent work unless a higher-level operation explicitly defines another recovery semantic.

Per-submission waiting is the narrow completion mechanism. Context-wide or resource-scoped synchronization is an optional capability for cases that need a broader boundary. Provider adapters retain native synchronization mechanics: Metal fences/events/barriers, CUDA stream/event behavior, OpenCL events and queue ordering, and Direct3D 12 fences/resource barriers are not interchangeable details. Their providers map generic dependencies onto native mechanisms without weakening correctness \cite{metal-sync,opencl-spec,d3d12-queues,d3d12-barriers,cuda-guide}.

This contract does not promise preemption. Cancellation is best-effort: the provider reports whether work was stopped before execution, cancellation was requested after work may have begun, or cancellation was too late. Resources remain retained until completion is observed. A device-loss or native execution failure must produce a failed completion and release retained resources only when the provider knows they are no longer in use.

The image document's mutation policy also needs a concurrency rule. Concurrent reads may share immutable resource snapshots. Concurrent writes to the same logical image revision must be serialized, rejected, or resolved through an explicitly defined conflict policy. The initial implementation SHOULD serialize commits per image and use revision checks to reject stale replacement attempts. The precise Swift synchronization mechanism (actor, lock-protected handle, or another approach) remains open.

\subsection{Execution completion and public asynchrony}

GPU and external-runtime work is commonly asynchronous. The public API should not imply that all provider operations are synchronous. An async operation application is a reasonable Swift-facing design: it awaits successful execution and validation before returning a completed logical result or committing replacement.

Internal implementations may still build a deferred representation and fuse work before materialization. That optimization must not change the public guarantee: when the API reports successful completion, the returned image state is valid and its resources remain alive. Cancellation and failure must release temporary resources safely and must not commit a partial replacement.

\section{Numerical and image-semantic contracts}

\subsection{Why names are insufficient}

An API that exposes \texttt{colorCorrection}, \texttt{blend}, or \texttt{resize} without specifying semantics permits backend implementations that disagree while claiming to be equivalent. Numerical and representational rules must be stated at the operation level.

\subsection{Required semantic dimensions}

Depending on the operation, the contract should define:
\begin{itemize}
  \item input/output data types and permitted precision;
  \item accumulation precision and rounding behavior;
  \item numerical range, clipping, saturation, and normalization;
  \item color space and transfer function used for arithmetic;
  \item straight versus premultiplied alpha treatment;
  \item channel order and channel-specific behavior;
  \item border/edge extension and kernel discretization;
  \item interpolation method and coordinate convention;
  \item deterministic behavior or allowed nondeterminism;
  \item an output error bound or comparison method where bitwise identity is not realistic.
\end{itemize}

For blending, the distinction between linear-light and transfer-encoded arithmetic can change the result materially. For image resizing, pixel-center conventions and border extension matter. For label maps, nearest-neighbor or a label-aware rule is required; bilinear interpolation can create invalid category values. For masks, binary and continuous-valued semantics must not be conflated.

\subsection{Floating-point requirements}

Floating-point support is a first-class data and operation concern. The API must distinguish storage precision from computation and accumulation precision. A resource stored as FP16 might be processed with FP32 accumulation, while another provider may only support a lower precision path. Qualification must compare these facts with the operation's required numerical contract.

The project should not claim universal FP64 support across GPUs. If FP64 is required by an operation or caller, providers lacking a qualified FP64 implementation must be rejected or an explicitly permitted, scientifically justified alternative must be selected.

\subsection{Scientific data preservation}

Display rendering and quantitative processing must remain distinct. A display transform (contrast, gamma, white balance, lookup table, or annotation overlay) must not silently overwrite original sensor intensities or measurement data. Acquisition layers should be preserved as source data; derived display and analysis layers should be separate logical outputs unless the caller explicitly requests replacement.

This distinction is especially important for microscopy and other scientific imaging, where numeric intensity, calibration, spatial scale, and label identity may carry measurement meaning.

\section{Errors, diagnostics, and observability}

The library should expose structured failures rather than backend-specific error strings alone. Useful categories include:
\begin{itemize}
  \item invalid operation parameters;
  \item incompatible input/layer type;
  \item unsupported format, precision, or semantic mode;
  \item no qualified implementation;
  \item provider/device unavailable or lost;
  \item resource allocation or import failure;
  \item synchronization or execution failure;
  \item output postcondition failure;
  \item cancellation;
  \item replacement revision conflict.
\end{itemize}

Diagnostics should identify the operation and semantic version, input descriptor, candidate provider/device, implementation version, rejection reason, and any attempted preparation step. Sensitive pointers, driver internals, or large pixel buffers should not be exposed as diagnostic payloads.

A fallback decision should be observable. The caller or log should be able to determine which implementation was selected and why, including provider/device, conversions, precision, and materialization points. This is necessary for reproducibility and performance analysis.

\section{Illustrative API shape}

The following sketches communicate the intended separation. They are not final Swift declarations and do not settle generic type erasure, ABI, actor, or concurrency details.

\subsection{Operation object and composition}

\begin{lstlisting}
let blur = GaussianBlur(
    sigma: 1.2,
    border: .reflect,
    precision: .fp32
)

let color = ColorCorrection(
    exposure: 0.25,
    contrast: 1.1,
    workingSpace: .linearSRGB
)

let pipeline = color.then(blur)
\end{lstlisting}

\subsection{Target and result policy}

\begin{lstlisting}
// New image identity; original document is unchanged.
let derived = try await image.applying(
    pipeline,
    to: .layer(originalLayer.id)
)

// Same logical image identity; commit a new revision atomically.
try await image.apply(
    pipeline,
    to: .layer(originalLayer.id)
)
\end{lstlisting}

\subsection{Provider and implementation registration}

\begin{lstlisting}
// The application selects the provider packages it deploys.
registry.registerProvider(CPUProvider())
registry.registerProvider(MetalProvider())

// Independent operation implementations target the same operation.
registry.register(CPUGaussianBlur())
registry.register(MetalGaussianBlur())
registry.register(OpenCVGaussianBlur())
\end{lstlisting}

The CPU/Metal/OpenCV implementation packages are not dependencies of the core protocol/runtime. A consuming application imports and registers the ones it chooses. Registration does not imply universal qualification: the common resolver evaluates each implementation specification against the current environment and pipeline policy.

\subsection{Qualification and resolution}

\begin{lstlisting}
// Conceptual internal runtime path; ordinary callers need not manage candidates.
let resolution = try runtime.resolve(
    pipeline: pipeline,
    input: source,
    environment: executionEnvironment.snapshot()
)

let output = try await runtime.execute(resolution.plan)
\end{lstlisting}

The runtime's inspection API may expose ranked candidates and structured reasons for diagnostics. The standard apply API should perform resolution internally. The application states semantic requirements and optimization preferences; it does not select CUDA or Metal merely because it knows the operating-system name.

\section{Testing and verification strategy}

The architecture is only useful if its contracts can be falsified. Tests must be organized around semantic invariants first and backend performance second.

\subsection{Operation conformance tests}

Each operation should have a backend-independent reference specification and a conformance suite. Tests should cover normal, boundary, degenerate, and invalid inputs. For floating-point operations, use declared tolerances and test metrics rather than assuming bitwise equality across all devices.

\subsection{Image and resource model tests}

Tests MUST cover:
\begin{itemize}
  \item shallow image/document derivation and immutable resource sharing;
  \item copy-on-write behavior and proof that an input is not changed by a pure operation;
  \item image and layer identity behavior across new-image and replacement modes;
  \item generic resource request validation: kind, capacity, alignment, memory domain, and owning target;
  \item imported-resource retention and source lifetime while the import remains alive;
  \item declared resource-access retention until submitted work completes;
  \item dependency ordering, typed submission output, failure propagation, and cancellation outcomes;
  \item context/resource synchronization boundaries where the provider advertises those capabilities;
  \item row stride, plane stride, alignment, ROI/view offsets, and non-contiguous data at the image-domain layer;
  \item external-resource lifetime and safe release after asynchronous execution;
  \item preservation of unrelated layers when a targeted layer is transformed;
  \item explicit composite rendering and no accidental flattening;
  \item failure atomicity when processing or output validation fails.
\end{itemize}

\subsection{Qualification tests}

For every operation/provider pair, test positive and negative capability cases. A qualification test must verify not only that a candidate is accepted, but that it is rejected for the correct reason when a required format, precision, border mode, dimension, or device feature is absent. Platform presence alone must never be used as proof of qualification.

\subsection{Cross-backend numerical equivalence}

For each operation, use fixed source fixtures and a trusted reference implementation. Compare output descriptors and numerical results. The test specification must record:
\begin{itemize}
  \item exact semantic parameters and input interpretation;
  \item absolute and relative error limits, or domain-specific metrics;
  \item treatment of NaN, infinity, signed zero, and out-of-range values where applicable;
  \item deterministic/repeatability requirements;
  \item whether the comparison is pixel-wise, region-wise, or task-level.
\end{itemize}

A backend is not conformant merely because the image \textquotedblleft{}looks similar\textquotedblright{}. Conversely, bitwise equality should not be required when floating-point order, precision, and hardware math legitimately differ within a documented error contract.

\subsection{Execution and performance tests}

Benchmark both individual operations and complete sequences. Record warm-up, compilation, synchronization, transfer, allocation, and execution costs separately. Measure end-to-end latency and throughput, memory pressure, and tail latency. Compare same-device resident pipelines with cross-provider transitions. A nominally faster kernel is not a system improvement if conversion and transfer costs dominate.

\subsection{CI and development constraints}

The design should support headless command-line builds and tests; it must not require an IDE for ordinary validation. Hardware-specific tests should be labeled and run only on matching CI agents. CPU reference and contract tests should run on every supported CI platform. A passing CPU-only test suite does not establish GPU correctness; provider-specific suites remain necessary.

\section{Implementation sequence}

The core library and its external implementations are separate delivery tracks. Core phases establish stable contracts, specifications, policy validation, environment observation, and the common resolver/runtime. Concrete provider and operation-implementation packages then apply those interfaces independently; they are not sequential phases of core-library development.

\begin{longtable}{@{}L{0.085\textwidth}L{0.24\textwidth}Z{0.60\textwidth}@{}}
\toprule
\textbf{Stage} & \textbf{Deliverable} & \textbf{Acceptance evidence} \\
\midrule
0 & Architecture and policy decisions & Stable ownership/dependency rules; specification model; precision semantics; policy scopes/composition; environment-state limits; package boundaries. \\
1 & Operation and implementation contracts & Immutable operation objects, operation specifications, one generic implementation protocol, declarative implementation specifications, stable identifiers/versions. \\
2 & General execution contracts & Provider/device/resource contracts, static capabilities, dynamic state observations, freshness and unknown-value semantics. \\
3 & Image-processing protocol implementation & Image/layer/operation/composition types, typed policy builder, normalized immutable policy, type-erased internal registry descriptors. \\
4 & Environment discovery and implementation registry & Provider discovery, implementation registration, snapshot identity, diagnostics, fake-provider/fake-implementation contract tests. \\
5 & Common qualification and resolver & Spec-driven eligibility, hard requirements, preferences, cost ranking, structured conflicts and unsatisfied-requirement diagnostics. \\
6 & Execution coordinator and pipeline planning & Resource reservation/revalidation, whole-pipeline or segment routing, execution, transactional result publication, cancellation/failure behavior. \\
7 & Independent implementation packs and reference application & At least one external provider pack and multiple operation implementation classes for the same operation; the sample application registers packages without changing core code. Other provider/algorithm packs can be developed independently. \\
8 & Measured optimization & Residency-aware planning, cost calibration, memory reuse, batching/fusion only with conformance evidence, and full-pipeline benchmarks. \\
9 & Production readiness & Versioned packages, platform/provider/operation support matrix, diagnostics, API docs, packaging, compatibility, release process and tested consumer application. \\
\bottomrule
\end{longtable}

Provider and algorithm implementation work is an independent workstream, not a technology checklist. CPU, Metal, CUDA, OpenCL, Direct3D 12, Vulkan, Accelerate, Core Image, and OpenCV packages are prioritized by actual target requirements, hardware access, maintainability, and conformance evidence. A package may be developed in parallel once its required contracts are stable. No one provider is required merely to complete a numbered core phase.

\section{Risks and decisions still open}

The current status of each foundational decision is tracked in the \href{https://github.com/Shennagyp/ImageProcessing/blob/main/docs/research/architecture/decisions/Decision-Register.md}{Architecture Decision Register}. Items marked proposed or deferred below are not silently treated as frozen public API.

\subsection{Decisions that must be resolved before freezing public API}

\begin{enumerate}
  \item \textbf{Image identity and concurrency:} whether the public handle is a Swift reference type, value wrapper around a shared document state, actor, or another model. The required semantics are stable identity, controlled mutation, snapshots, and safe asynchronous commits.
  \item \textbf{Layer identity policy:} whether layer IDs are globally unique or scoped to an image/document. This must be decided before serialization or long-lived external references are added.
  \item \textbf{Resource handle ABI:} the boundary between Swift, C/C++, and provider adapters. A narrow C ABI may be useful for stable cross-language interoperability, but it must not leak native implementation handles into the semantic API.
  \item \textbf{Cross-platform language surface:} a Swift-first API is natural for Apple platforms, but a universal Swift application/runtime surface across every target operating system must be validated rather than assumed. The semantic core and provider ABI should be designed independently of a promise that Swift tooling is identical on all platforms.
  \item \textbf{Lazy versus eager result contract:} internal deferred execution is desirable, but the public completion/materialization boundary must be explicit. The initial async apply API should return only after its promised output state is valid.
  \item \textbf{Color-management defaults:} default working space, scientific non-color data behavior, alpha semantics, and conversion policy need a dedicated specification.
  \item \textbf{Operation contract representation:} determine which properties are compile-time generic types and which are runtime descriptors. Do not force all dynamic image formats into unwieldy generic type parameters.
  \item \textbf{Provider plug-in lifecycle:} define loading, version compatibility, failure isolation, and whether providers are statically linked, dynamically loaded, or both.
\end{enumerate}

\subsection{Risks and mitigations}

\begin{longtable}{@{}L{0.25\textwidth}Z{0.69\textwidth}@{}}
\toprule
\textbf{Risk} & \textbf{Mitigation} \\
\midrule
Lowest-common-denominator provider API becomes too weak & Keep the common protocol small and add optional capability protocols; preserve native handles inside provider adapters. \\
Provider interface becomes a giant abstraction & Separate provider lifecycle/resource execution from operation implementations; register capabilities at operation granularity. \\
Silent numerical drift between implementations & Version semantic contracts, define numerical guarantees, and require conformance tests for every implementation. \\
Greedy routing causes repeated transfers & Track residency and transition costs over composed sequences; benchmark whole pipelines. \\
Unclear copy/replace behavior & Separate logical document identity from storage identity; transactional replacement; aliasing tests. \\
Layer model becomes an over-general object hierarchy & Use a small shared layer identity/property contract with typed content kinds; add a new layer kind only when it has distinct semantics. \\
External framework conflicts with runtime ownership & Treat Core Image, OpenCV, ONNX Runtime, and similar systems as implementation adapters with explicit control/residency boundaries. \\
Swift portability is assumed without evidence & Keep core semantics language-neutral and verify target toolchains and packaging before promising universal Swift support. \\
\bottomrule
\end{longtable}

\section{Consolidated architecture rules}

The following rules are the compact invariants that future API and implementation decisions must preserve.
\begin{enumerate}
  \item An image is a logical layered object, not a raw pixel allocation.
  \item A layer has a semantic content kind; raster, mask, label, and annotation data are not interchangeable.
  \item Pixel storage is a resource with explicit descriptor, ownership, lifetime, residency, and synchronization.
  \item An operation is an immutable object describing meaning and parameters, not a backend-specific executable.
  \item Operation semantics are versioned independently from implementation versions.
  \item One generic \texttt{ImageOperationImplementation} protocol is specialized by an associated operation type; multiple independent classes may implement the same operation.
  \item Operation specifications describe semantic requirements and permitted variation. Implementation specifications declare one implementation's static, tested support and provider/resource requirements.
  \item Implementation classes do not each implement bespoke qualification logic. The common resolver evaluates declarative specifications.
  \item \texttt{ImageProcessing\allowbreak Protocol} owns image semantics, operation and implementation contracts, declarative policy types, and qualification result vocabulary.
  \item \texttt{GeneralExecution\allowbreak Protocol} is a lower-level, image-agnostic dependency describing providers, devices, resources, submission, completion, capabilities, and observable execution state; it MUST NOT depend on image-processing concepts.
  \item Provider packages implement the generic execution contract. Operation implementation packages implement image-processing contracts. The consuming application chooses and registers both.
  \item Platform is an environment fact, not a default provider preference. Eligibility follows actual provider discovery, declared implementation requirements, and device capabilities.
  \item The execution environment distinguishes static capability from dynamic availability/load and records telemetry freshness; unknown load is not treated as zero.
  \item One common resolver evaluates operation semantics, implementation specifications, current environment state, and pipeline policy. It filters invalid candidates and ranks only eligible candidates.
  \item Requirements determine admissibility; preferences guide ranking and can never override semantic or hard execution requirements.
  \item Pipeline qualification is declarative and fluent. It produces immutable policy values rather than injected per-implementation code.
  \item Hard policy requirements compose conjunctively and are validated when built/composed. Empty intersections and contradictory assignments produce provenance-rich diagnostics.
  \item Soft preferences use explicit objective dimensions and declared trade-off semantics; they are not intersected as sets.
  \item Dynamic qualification is snapshot-based. Resource reservation and critical state are revalidated before execution; stale plans may be resolved again only within original hard constraints.
  \item Routing considers queue delay, transfer, conversion, allocation, compute, synchronization, materialization, residency, and pipeline transitions—not kernel time alone.
  \item CPU fallback is permitted only when it satisfies all hard requirements and the pipeline does not require another provider.
  \item New-image and replace-image APIs express logical behavior; neither promises in-place memory mutation.
  \item Replacement is transactional and preserves the old image state on failure.
  \item Operation targets (layer, composite, or explicit multi-input set) are never inferred in a destructive or ambiguous way.
  \item Numerical meaning is part of the API contract and is tested across implementations.
  \item No node editor or public visual graph model is included in this design.
\end{enumerate}

\section{Conclusion}

The design should not imitate one existing framework wholesale. OpenCV contributes mature image algorithms and practical array/storage lessons; PyTorch demonstrates a common numerical object and dispatch architecture; ONNX Runtime demonstrates capability-based execution-provider partitioning; Core Image demonstrates useful deferred image evaluation; and native compute APIs demonstrate why resources, command submission, and synchronization require provider-specific implementations.

The proposed unification is deliberately narrower: a layered image object model, immutable operation objects, typed composition, explicit new-versus-replace semantics, the domain-facing \texttt{ImageProcessing\allowbreak Protocol}, and the lower-level, image-agnostic \texttt{GeneralExecution\allowbreak Protocol}. Provider packages implement the general execution contract; independent operation-implementation classes implement the image-domain contract. The runtime resolves concrete implementation/provider/device bindings using declarative specifications, current environment observations, and pipeline policy, then coordinates execution without collapsing those responsibilities.

The most consequential boundary is the separation between \emph{what an operation means} and \emph{which implementation executes it}. The most consequential data boundary is the separation between \emph{logical image/layer identity} and \emph{physical pixel storage}. Preserving both boundaries is the foundation for portability, numerical consistency, safe mutation, and efficient heterogeneous execution.

\clearpage
\appendix
\section{Glossary}

\begin{longtable}{@{}L{0.25\textwidth}Z{0.69\textwidth}@{}}
\toprule
\textbf{Term} & \textbf{Meaning in this document} \\
\midrule
Image document & Logical image object consisting of a canvas, ordered layers, and metadata. \\
Layer & Typed semantic component of an image document. \\
Pixel resource & Opaque reference to numerical pixel/array storage. \\
Operation object & Immutable parameterized description of a computation and its contract. \\
Operation implementation & Concrete algorithm class implementing one operation; it publishes a declarative implementation specification and executes through a provider or managed runtime. \\ 
Image operation-implementation contract & Generic typed protocol specialized by an associated operation type. \\ 
Execution candidate & Concrete binding of an implementation instance to a provider/runtime and device/execution target, evaluated against an environment snapshot. \\
Compute provider & Execution substrate adapter exposing devices, resources, submission, completion, and synchronization. \\
Operation specification & Declarative statement of an operation's semantic requirements and permitted behavior. \\
Implementation specification & Declarative, versioned statement of one implementation's tested support, constraints, provider requirements, and resource characteristics. \\
Execution environment snapshot & Time-bounded view of available providers/devices, capabilities, resource state, and observable workload/availability. \\
Qualification policy & Immutable declarative requirements and preferences supplied to the common resolver. \\
Qualification/resolution & Shared evaluation that filters ineligible implementations and ranks eligible candidates against policy and current environment state. \\
Requirement & Hard condition an implementation must satisfy to remain eligible. \\
Preference & Soft objective used to rank eligible candidates; it cannot relax a requirement. \\


Residency & Current location/access domain of a resource (host, device, shared, external, or managed). \\
Materialization & Evaluation of a deferred representation into a concrete resource/output. \\
Logical replacement & Atomic update of an image document's state while retaining its logical identity. \\
Physical in-place execution & Reuse/overwrite of an existing allocation; an optimization independent of logical replacement. \\
\bottomrule
\end{longtable}

\section{Reference package structure (provisional)}

The package names below express dependency boundaries, not a requirement to create all packages in the first milestone.
\begin{lstlisting}
GeneralExecutionProtocol/
  Provider, Device, Context, Resource, Allocation/Import
  Work, Submission, Completion, Cancellation, Synchronization
  CapabilityDescriptor, Availability/StateObservation

ImageProcessingProtocol/
  Image, Layer, PixelResourceDescriptor
  ImageFunction, Operation, OperationSpecification
  ImageOperationImplementation
  ImplementationSpecification, QualificationPolicy

ImageProcessingRuntime/
  Registry, EnvironmentDiscovery, Snapshot
  CommonQualifierResolver, CostModel, ExecutionPlanner
  TypeErasure, ResourceReservation, ExecutionCoordinator
  ImageResult, ReplacementTransaction

ExternalExecutionProviderPacks/
  CPU, Metal, CUDA, OpenCL, Direct3D12, Vulkan

ExternalOperationImplementationPacks/
  GaussianBlur: CPU, Metal, CUDA, OpenCV, ...
  Resize: CPU, Metal, CUDA, OpenCV, ...
  CoreImage, Accelerate, OpenCV, ONNXRuntime, CoreML adapters
\end{lstlisting}

Dependency rules:
\begin{itemize}
  \item \texttt{GeneralExecution\allowbreak Protocol} has no dependency on image-processing types.
  \item \texttt{ImageProcessing\allowbreak Protocol} may depend on \texttt{GeneralExecution\allowbreak Protocol} for generic execution contracts, but the reverse dependency is forbidden.
  \item Provider implementations conform to \texttt{GeneralExecution\allowbreak Protocol}.
  \item Operation adapters conform to \texttt{ImageProcessing\allowbreak Protocol}; native operation adapters may also consume a concrete provider implementation.
  \item \texttt{ImageProcessingRuntime} integrates the protocols and implements registry/selection/execution orchestration without becoming a provider.
\end{itemize}

\section{Reference API vocabulary (provisional)}

The names below are vocabulary candidates, not frozen declarations.
\begin{longtable}{@{}L{0.28\textwidth}Z{0.66\textwidth}@{}}
\toprule
\textbf{Name} & \textbf{Intended meaning} \\
\midrule
\texttt{Image} & Logical image document/handle. \\
\texttt{Layer} & Typed content entry in an image document. \\
\texttt{PixelResource} & Opaque numerical storage handle. \\
\texttt{ImageFunction} & Typed, composable function over image-domain inputs and outputs. \\
\texttt{Operation} & Immutable semantic operation description/instance. \\
\texttt{OperationSpecification} & Semantic requirements and permitted behavior for one operation. \\
\texttt{ImageOperationImplementation} & Generic typed contract for one implementation class, specialized by its associated operation type. \\
\texttt{ImplementationSpecification} & Declarative static support and execution requirements of one implementation class. \\
\texttt{ComputeProvider} & General execution technology adapter exposing generic execution lifecycle and resources. \\
\texttt{QualificationPolicy} & Immutable requirements and preferences attached to a function/pipeline through the fluent API. \\
\texttt{ExecutionEnvironmentSnapshot} & Time-bounded provider/device capability and runtime-state observations. \\
\texttt{ResolutionResult} & Candidate eligibility, ranking, selected plan, cost assumptions, and diagnostics. \\
\texttt{ExecutionPlan} & Internal plan for a composed operation/function. \\
\texttt{applying} & Produce a new logical image result. \\
\texttt{apply} & Replace image state transactionally. \\
\bottomrule
\end{longtable}

\section{References}
\begin{thebibliography}{99}

\bibitem{opencv-mat}
OpenCV Contributors,
\emph{cv::Mat Class Reference}, OpenCV documentation, 4.x series.
\newblock \url{https://docs.opencv.org/4.x/d3/d63/classcv_1_1Mat.html}
\newblock Accessed 20 November 2026.

\bibitem{opencv-gapi}
OpenCV Contributors,
\emph{Graph API: High-Level Design Overview}, OpenCV documentation, 4.12 series.
\newblock \url{https://docs.opencv.org/4.12.0/de/d4d/gapi_hld.html}
\newblock Accessed 20 November 2026.

\bibitem{pytorch-tensor}
PyTorch Contributors,
\emph{Tensor Attributes}, PyTorch documentation.
\newblock \url{https://docs.pytorch.org/docs/stable/tensor_attributes.html}
\newblock Accessed 20 November 2026.

\bibitem{pytorch-dispatch}
PyTorch Contributors,
\emph{Extending PyTorch: Dispatcher and Native Operator API}, PyTorch documentation.
\newblock \url{https://docs.pytorch.org/docs/stable/notes/extending.html}
\newblock Accessed 20 November 2026.

\bibitem{pytorch-custom-backend}
PyTorch Contributors,
\emph{Extending the Dispatcher for a New Backend in C++}, PyTorch tutorials.
\newblock \url{https://docs.pytorch.org/tutorials/advanced/extend_dispatcher.html}
\newblock Accessed 20 November 2026.

\bibitem{onnx-ir}
ONNX Contributors,
\emph{ONNX Intermediate Representation Specification}, ONNX documentation.
\newblock \url{https://onnx.ai/onnx/repo-docs/IR.html}
\newblock Accessed 20 November 2026.

\bibitem{ort-architecture}
ONNX Runtime Contributors,
\emph{ONNX Runtime High-Level Design}, ONNX Runtime documentation.
\newblock \url{https://onnxruntime.ai/docs/reference/high-level-design.html}
\newblock Accessed 20 November 2026.

\bibitem{ort-providers}
ONNX Runtime Contributors,
\emph{Execution Providers}, ONNX Runtime documentation.
\newblock \url{https://onnxruntime.ai/docs/execution-providers/}
\newblock Accessed 20 November 2026.

\bibitem{coreimage-ciimage}
Apple Inc.,
\emph{CIImage}, Core Image documentation.
\newblock \url{https://developer.apple.com/documentation/coreimage/ciimage}
\newblock Accessed 20 November 2026.

\bibitem{coreimage-cicontext}
Apple Inc.,
\emph{CIContext}, Core Image documentation.
\newblock \url{https://developer.apple.com/documentation/coreimage/cicontext}
\newblock Accessed 20 November 2026.

\bibitem{coreimage-overview}
Apple Inc.,
\emph{Core Image}, Apple Developer Documentation.
\newblock \url{https://developer.apple.com/documentation/coreimage}
\newblock Accessed 20 November 2026.

\bibitem{metal-sync}
Apple Inc.,
\emph{Resource Synchronization}, Metal documentation.
\newblock \url{https://developer.apple.com/documentation/metal/resource-synchronization}
\newblock Accessed 20 November 2026.

\bibitem{cuda-guide}
NVIDIA Corporation,
\emph{CUDA Programming Guide}, NVIDIA documentation.
\newblock \url{https://docs.nvidia.com/cuda/cuda-programming-guide/}
\newblock Accessed 20 November 2026.

\bibitem{opencl-spec}
The Khronos Group,
\emph{The OpenCL Specification}, unified OpenCL API specification, version 3.1.2, 17 September 2026.
\newblock \url{https://registry.khronos.org/OpenCL/specs/unified/html/OpenCL_API.html}
\newblock Accessed 20 November 2026.

\bibitem{d3d12-queues}
Microsoft,
\emph{Command Queues and Command Lists}, Direct3D 12 documentation.
\newblock \url{https://learn.microsoft.com/en-us/windows/win32/direct3d12/command-queues-and-command-lists}
\newblock Accessed 20 November 2026.

\bibitem{d3d12-barriers}
Microsoft,
\emph{Important Changes from Direct3D 11 to Direct3D 12}, Direct3D 12 documentation.
\newblock \url{https://learn.microsoft.com/en-us/windows/win32/direct3d12/important-changes-from-directx-11-to-directx-12}
\newblock Accessed 20 November 2026.

\end{thebibliography}

\end{document}
